\documentclass[11pt]{ctexart}
\usepackage[a4paper,margin=1in]{geometry}
\usepackage{graphicx}
\usepackage{xcolor}
\usepackage{hyperref}
\definecolor{refgray}{gray}{0.45}
\title{\textbf{市场最新观点汇总}\\[0.4em]{\large 盈利周期加速与供给侧再定价成为主线，AI资本开支的结构性支撑正在重塑宏观与行业格局}}
\author{KC桌面 自动生成}
\date{260514}
\begin{document}
\maketitle
\section*{一页摘要}
\begin{itemize}
  \item 美股盈利周期正在加速，标普500成分股1Q26盈利超预期中位数创四年新高，滚动式复苏的早期阶段已确认。
  \item AI资本开支正从科技行业内部故事转变为宏观经济的结构性支撑力量，对冲了能源冲击对消费的拖累。
  \item 亚洲工业超级周期启动，固定投资增速预计三倍于此前水平，工业周期强度系统性超越能源冲击。
  \item 中国电动车市场定价权来源从规模切换为产品迭代速度，改款车型成为价格新锚点，行业竞争逻辑转变。
  \item 欧洲能源巨头1Q26业绩超预期来自下游和交易业务，管理层保守的资本分配暗示未来股东回报上调空间。
\end{itemize}
\section{宏观与利率：盈利周期加速，AI资本开支成为结构性支撑}
\textbf{市场真正低估的不是风险，而是盈利周期正在加速；AI资本开支正在从科技行业内部故事转变为宏观经济的结构性支撑力量。}

\begin{itemize}
  \item 标普500成分股2026年第一季度盈利超预期幅度的中位数达到6\%，创下四年以来最强水平，盈利修正广度从5\%迅速回升至22\%，显示企业盈利韧性超预期。
  \item 美国经济正经历从消费驱动向AI投资驱动的范式切换，非住宅固定投资2026年预计增长7\%，2027年加速至8\%，科技巨头资本支出2027年有望突破1万亿美元。
  \item 韩国经济正从K型复苏走向全面反弹，2026年一季度GDP环比增长1.7\%为2020年以来最强，私人消费、建筑投资和资本开支全面回暖，增长分化正在弥合。
  \item 韩国央行可能从2026年四季度开始加息，终端利率上看3.25\%，核心CPI可能升至两年最高水平，财政空间比市场预期更大。
  \item 美国实际消费增速受油价冲击将从2025年的2.1\%降至2026年的1.8\%，但AI投资驱动的资本支出足以对冲消费拖累，GDP增速仍可维持2.3\%左右。
\end{itemize}
\begin{figure}[htbp]
\centering
\includegraphics[width=0.88\linewidth]{figures/fig_018.jpg}
\caption{Exhibit 1 - \textbackslash{}- Sector/style recommendations: We see cyclicals outperforming defensives, which is consistent with our rolling recovery/earnings acceleration thesis. Small cap earnings are infle}
\end{figure}
\begin{figure}[htbp]
\centering
\includegraphics[width=0.88\linewidth]{figures/fig_020.jpg}
\caption{Exhibit 1 - Exhibit 1: US Economics Year-Ahead Outlook: Forecast Summary <table><tr><td colspan="5">4Q/4Q \% change, except where noted</td></tr><tr><td></td><td>2024A</td><td>2025A</td><td>202}
\end{figure}
{\small\color{refgray}\textbf{References:} [R011] 市场真正低估的不是风险，而是盈利周期正在“跑热”; [R012] 市场真正低估的不是消费，而是AI资本开支的结构性替代效应; [R010] 韩国经济正在从“K型复苏”走向“全面反弹”，市场低估了消费和财政的接力效应}

\section{AI/算力：供给侧再定价与芯片供应拐点}
\textbf{市场低估了AI算力供给侧的再定价，芯片供应约束正在松动，数据中心交付节奏从等待芯片切换为交付产能。}

\begin{itemize}
  \item 某投行将2030年AI数据中心系统TAM从1.4万亿美元上调至1.7万亿美元，CAGR达45\%，核心逻辑是2026-2027年AI公司IPO和token经济学改善将验证资本支出可持续性。
  \item 乌兰察布数据中心四月新增装机125MW创历史新高，超过整个一季度合计71MW，芯片供应瓶颈实质性解除，积压订单开始集中交付。
  \item 半导体设备行业景气判断正从利润驱动转向订单驱动，Ulvac订单指引从增长24\%上调至37\%创历史新高，AI存储、OLED和稀土需求是三大驱动力。
  \item 制程迁移正在重塑IP公司价值曲线，eMemory受益于OLED驱动芯片向16nm和电源管理芯片向55nm迁移，单颗芯片IP价值阶梯式跃升。
  \item 砷化镓代工产业中，卫星通信正取代智能手机成为稳懋2026年主要增长引擎，预计贡献中双位数营收占比，而AI数据中心业务占比仍仅为中个位数。
\end{itemize}
\begin{figure}[htbp]
\centering
\includegraphics[width=0.88\linewidth]{figures/fig_022.jpg}
\caption{Figure 1 - Analyst navin.killa@UBS.com +852-2971 7594 Figure 1: Implied net adds in Ulanqab IDC total power (MW, IT power x PUE) <details> <summary>bar</summary> | Month | Implied MW net adds}
\end{figure}
\begin{figure}[htbp]
\centering
\includegraphics[width=0.88\linewidth]{figures/fig_032.jpg}
\caption{Exhibit 1 - stay UW: We believe expectations for WIN Semi's InP (Indium Phosphide) business are too bullish, as AI datacenter still only accounts for mid-single digits \% of revenue. \# WIN Semi}
\end{figure}
{\small\color{refgray}\textbf{References:} [R002] 市场真正低估的不是需求，而是AI算力供给侧的再定价; [R014] 乌兰察布数据中心四月新增装机创纪录：市场低估的不是需求，而是芯片供应拐点后的交付弹性; [R019] 半导体设备股的真正信号不在利润，而在订单的结构性重置; [R022] 市场低估了制程迁移带来的IP价值重估，而非仅仅AI需求; [R021] 市场真正低估的不是需求，而是供给侧的再定价}

\section{能源与大宗：欧洲油企重建分配能力，工业超级周期启动}
\textbf{欧洲油企在动荡中展现出盈利结构变化，下游和交易业务驱动超预期，管理层保守的资本分配暗示未来股东回报上调空间。}

\begin{itemize}
  \item 欧洲五大油企1Q26业绩超预期主要来自下游和交易业务，BP客户与产品板块利润超预期30\%，Shell化工与产品板块超预期35\%，上游和LNG业务反而普遍低于预期。
  \item 上游涨价红利将在2Q26显现，合约价格滞后效应导致1Q26上游实现价格比现货低约5美元/桶，若油价维持高位，2Q26上游实现价格环比可能改善约40\%。
  \item 管理层在资本分配上异常保守，优先修复资产负债表而非加速分红，但当前分红和回购水平远未反映真实现金流生成能力，一旦环境稳定分配力度上调是大概率事件。
  \item 亚洲工业超级周期正在启动，固定资本投资预计从当前约11万亿美元增长至2030年的16万亿美元，年复合增长率7\%，约为2023-2025年增速的三倍。
  \item 工业周期强度系统性超越能源冲击，亚洲制造业PMI回升，工业产出创四年新高，非科技类出口持续回暖，资本品进口显示投资周期走强。
\end{itemize}
\begin{figure}[htbp]
\centering
\includegraphics[width=0.88\linewidth]{figures/fig_010.jpg}
\caption{Exhibit 1 - Earnings beat on elevated expectations, driven primarily by downstream margins and trading/optimisation gains. 1Q upstream pricing remained muted by contractual lags; stronger real}
\end{figure}
\begin{figure}[htbp]
\centering
\includegraphics[width=0.88\linewidth]{figures/fig_005.jpg}
\caption{Exhibit 2 - \# The strength of the industrial cycle is outweighing the energy shock At the centre of two countervailing forces: Our recent research has been focused on two issues – 1) the stren}
\end{figure}
{\small\color{refgray}\textbf{References:} [R007] 市场真正低估的不是地缘溢价，而是欧洲油企正在重建的分配能力; [R003] 市场真正低估的不是能源风险，而是工业超级周期的广度}

\section{地产与消费：电动车定价回暖，宠物食品品牌分化}
\textbf{中国电动车定价权从规模切换为产品迭代速度，宠物食品行业进入品牌资产驱动阶段，消费板块整体仍承压但结构性机会显现。}

\begin{itemize}
  \item 4月多家头部电动车品牌加权平均零售价环比改善，改款车型是核心驱动力，蔚来5566系列改款后每款车涨价约10\%，吉利银河改款后涨幅超10\%。
  \item 五一后电动车周订单普遍回落，但特斯拉、蔚来、理想相对抗跌，核心分界线是产品周期位置而非价格战力度，市场正从谁降价谁拿订单转向谁有新车谁有话语权。
  \item 宠物食品行业4月线上GMV仅增长1.1\%，但头部企业中国宠物食品同比增长29.2\%，猫主粮韧性远强于狗粮，抖音和京东渠道增长远优于天猫。
  \item A股1Q26盈利底正在确认，按公司数量计算的净miss率从-23\%收窄至-12.5\%，加权surprise转正至+0.3\%，但结构分化严重：工业和金融率先走出底部，下游消费仍在泥潭。
  \item 工业板块改善核心动力来自资本品子领域（AIDC设备、电池设备、自动化与机器人），领先企业展现定价能力能通过提价部分抵消成本上涨。
\end{itemize}
\begin{figure}[htbp]
\centering
\includegraphics[width=0.88\linewidth]{figures/fig_007.jpg}
\caption{Exhibit 1 - XPeng +0.6\% MoM: Retail price of MONA M03 increased by Rmb6.5k following the facelift, while other models remained flat. Further uplifts will likely depend on the launch of GX. NIO}
\end{figure}
\begin{figure}[htbp]
\centering
\includegraphics[width=0.88\linewidth]{figures/fig_008.jpg}
\caption{Exhibit 1 - \# Vicky Wu Equity Strategist Vicky.Wu@morganstanley.com +852 3963-3928 Exhibit 1: History of Chinese equity market quarterly earnings surprises by number of companies (MSCI China a}
\end{figure}
{\small\color{refgray}\textbf{References:} [R004] 4月电动车定价回暖：改款周期正在重塑价格战逻辑，而非终结它; [R006] 五一后销量骤降不是新闻，真正关键的是“产品周期”正在重新划分阵营; [R015] 宠物食品行业真正的分化不在大盘增速，而在品牌势能能否转化为渠道议价权; [R005] A股盈利底正在确认，但真正的分化才刚刚开始}

\section{区域市场：韩国全面反弹，人民币中间价信号被忽略}
\textbf{韩国经济从K型复苏走向全面反弹，人民币中间价定价框架可能经历结构性调整，市场对汇率信号的认知不够充分。}

\begin{itemize}
  \item 韩国2026年GDP增速预计从2025年的1.0\%飙升至2.8\%，私人消费增速从1.3\%提升至2.2\%，高收入群体消费和旅游支出拉动内需，财政空间充裕可追加预算。
  \item 韩国通胀压力正从能源端向服务业传导，核心CPI可能升至2.3\%，央行加息路径比市场预期更鹰派，终端利率可能到3.25\%。
  \item 人民币中间价模型预测值降至6.7902，较前次低529个基点，但变化主要来自一篮子货币隔夜波动（卢布、墨西哥比索等），而非央行主动调整逆周期因子。
  \item 模型误差近期收窄至零附近，暗示市场对中间价预期更一致，政策信号更清晰，5月中美领导人会晤和7月底政治局会议或成汇率关键节点。
\end{itemize}
\begin{figure}[htbp]
\centering
\includegraphics[width=0.88\linewidth]{figures/fig_016.jpg}
\caption{Exhibit 1 - Chief Korea/Taiwan Economist Kathleen.Oh@morganstanley.com +852 2848-7340 Exhibit 1: We expect a broad-based rebound in GDP growth to 2.8\% YoY in 2026 vs. 1.0\% in 2025 <details> <s}
\end{figure}
\begin{figure}[htbp]
\centering
\includegraphics[width=0.88\linewidth]{figures/fig_017.jpg}
\caption{Exhibit 2 - Exhibit 2 : Korea - forecast summary <table><tr><td></td><td>2024</td><td>2025</td><td>2026F</td><td>2027F</td></tr><tr><td>Real GDP, \%Y</td><td>2.0</td><td>1.0</td><td>2.8</td><td}
\end{figure}
{\small\color{refgray}\textbf{References:} [R010] 韩国经济正在从“K型复苏”走向“全面反弹”，市场低估了消费和财政的接力效应; [R013] 人民币中间价模型正在发出一个被市场忽略的定价信号}

\section{风险偏好与对冲：泡沫风险指标逼近警戒，对冲策略需升级}
\textbf{市场真正低估的不是AI泡沫本身，而是对冲泡沫破裂的工具正在被重新定价，传统保护性看跌期权在单边急涨市场中失效。}

\begin{itemize}
  \item 某投行泡沫风险指标显示美国科技板块读数达0.7，纳斯达克0.65，韩国KOSPI高达1.0，日本日经0.85，已接近或超过历史上泡沫形成和见顶阶段的典型区间。
  \item 传统保护性看跌期权面临行权价选择和时机选择双重缺陷，在单边急涨市场中成本沉没或吞噬收益，基于波动率的对冲结构（如回溯看跌逻辑的衍生品）成为兼顾成本与保护的解决方案。
  \item 欧洲地缘风险溢价正在消退，亚洲AI股票中最差看涨期权展现出惊人性价比，市场对尾部风险的定价可能不够充分。
  \item 美股滚动式复苏的核心支撑是企业盈利韧性而非政策，标普500远期市盈率已从峰值压缩18\%，市场在过去六个月完成了大量定价工作。
\end{itemize}
\begin{figure}[htbp]
\centering
\includegraphics[width=0.88\linewidth]{figures/fig_001.jpg}
\caption{Exhibit 4 - Equities led stress lower as they posted the largest absolute change in stress for the fourth straight week (Exhibit 4). This came as the S\&P 500 recorded its sixth consecutive wee}
\end{figure}
\begin{figure}[htbp]
\centering
\includegraphics[width=0.88\linewidth]{figures/fig_002.jpg}
\caption{Exhibit 3 - </details> Source: BofA Global Research. \textbackslash{}*Latest as of 8-May-26. Disclaimer: The indicator identified above as BofA GFSI is intended to be an indicative metric only and may not be}
\end{figure}
{\small\color{refgray}\textbf{References:} [R001] 市场真正低估的不是AI泡沫本身，而是对冲泡沫破裂的工具正在被重新定价; [R011] 市场真正低估的不是风险，而是盈利周期正在“跑热”}

\section*{结语}
综合来看，当前市场主线是盈利周期加速与供给侧再定价的共振，AI资本开支的结构性支撑正在重塑宏观与行业格局，但各板块分化显著，投资者需关注产品周期、成本传导能力和现金流可预测性等微观变量。
\vfill
{\small\color{refgray}Personal reading notes and learning share only. Not investment advice.}
\end{document}
